% Generated by tex/build_swebok_structure.py from tex/chapters/ch01/03_ソフトウェア要求の基礎.tex
\subsubsection{技術制約}

\term{技術制約}とは、特定の自動化技術やインフラを使うこと、または使わないことを指定する要求である。ここでは「名前のある技術を指定しているか」に注目すると見分けやすい。

\par\smallskip\noindent\refstepcounter{table}\label{tab:ch01-04}\textbf{表 \thetable: 第01章 ソフトウェア要求 表04}\par\smallskip
\begin{tabularx}{0.97\linewidth}{@{}p{.25\linewidth}X@{}}
\toprule
対象 & 例 \\
\midrule
プラットフォーム & Windows、macOS、Android、iOS など。 \\
プログラミング言語 & Java、C++、C\#、Python など。 \\
ブラウザ & Chrome、Safari、Edge など。 \\
データベース & Oracle、SQL Server、MySQL など。 \\
禁止事項 & ポインタ利用禁止、特定暗号方式の禁止、外部クラウド利用禁止など。 \\
\bottomrule
\end{tabularx}

技術制約は、既存システムとの連携、組織標準、保守要員のスキル、セキュリティ方針、ライセンス方針などから生じることが多い。
